




\drghosh{\textbf{Title.} Title is good. I like it since it poses a clear puzzle ('escape' vs. 'trap'). Make sure the intro and conclusion explicitly answer which side wins. Right now the event-study suggests the 'trap' interpretation (negative coefficients), but this is never stated upfront.}





\drghosh{\textbf{Intro length \& structure.} Currently \textasciitilde{}1/2 page; for a top-field journal, aim for 3–4 pages. Standard AER/JDE structure: (1) motivation; (2) the research question; (3) institutional setting (Nepal civil war 1996–2006, Maoist insurgency, \textasciitilde{}17,000 deaths); (4) data + identification (one paragraph); (5) preview of findings with magnitudes; (6) contribution to $\geq$3 literatures (conflict and migration; childhood exposure to violence; long-run effects of civil war); (7) road map. The current text reads as three loose observations, not an introduction.}



\drghosh{\textbf{History placement.} The Lahore/Gurkha history is interesting and shows migration is structurally embedded in Nepalese life. But it belongs in a 'Background' section, not in the first paragraph of the intro. Lead with the puzzle, not the history.}

International migration of Nepalese people has been a common phenomenon even before nineteenth century when they used to travel to Lahore to join the army of the Sikh ruler Ranjit Singh which was further institutionalised with the recruitment of Nepalese to the British Gurkhas in 1815-16. 


\drghosh{\textbf{Broken sentence.} 'Seddon et al. (2002) Foreign migration has been a by-product of stagnant rural economy' — sentence is broken. Either: 'Seddon et al. (2002) argue that foreign migration is a by-product of a stagnant rural economy.' Or cite a specific finding. Right now this looks like an unfinished note.}

\cite{Seddon2002} Foreign migration has been a by-product of stagnant rural economy 



\drghosh{\textbf{Two-pathway theory.} You state the two-pathway theory (destroy vs. build social capital) but never connect it to your research question. This is actually a useful framing for your title's question: ex ante, conflict could push people out (escape) or trap them in (capital destruction). Use this paragraph to motivate why the sign of the long-run effect is theoretically ambiguous, then say which one you find.}



\drghosh{\textbf{Missing intro elements.} No mention of: (i) what you actually find; (ii) why Nepal is a useful case (one of the most migration-dependent economies in the world; \textasciitilde{}25\% of GDP from remittances; recent civil war with sharp geographic variation); (iii) what's novel relative to Bohra-Mishra \& Massey (2011) — they study contemporaneous migration during conflict, you study long-run effects on a cohort that was children during conflict; (iv) policy implications.}

\cite{Voors2012} advocate that there are two pathways for an economy after civil war; if it destroys the social capital or increases impatience then it will reduce investment and growth, however if it triggers institutional improvements or alters preferences towards saving then income might even be higher than before. 